%% kapa: exam=1 round=29 year=2018 period=1 form=A subject=real_estate subject_ko=부동산학원론 seq=13 number=93
%% source: Q-Net 감정평가사 기출문제 (cst00309) artlSeq=5202808 file=extracted/1008/5202808/2018년 제29회 감정평가사 1차시험 기출문제/2018년 제29회 감정평가사 1차시험 기출문제/1교시 A형.pdf
%% conversion: kordoc 4.24.51 -> tools/gen1.py (auto); review=proofread-sonnet; flags=none
\begin{question}{93}
부동산 투자타당성 분석기법에 관한 설명으로 옳지 않은 것은?
\begin{choices}
  \item 수익성지수는 투자개시시점에서의 순현가와 현금지출의 현재가치 비율이다.
  \item 내부수익률법은 화폐의 시간가치를 고려한다.
  \item 동일한 투자안에 대해서 복수의 내부수익률이 존재할 수 있다.
  \item 내부수익률은 순현가가 ‘0’이 되는 할인율이다.
  \item 순현가법에 적용되는 할인율은 요구수익률이다.
\end{choices}
\end{question}
\answer{1}
